% =====================================================================
%  Harald Høffding, Om Pessimismen
%  In: Vor Ungdom.  Tidsskrift for Opdragelse og Undervisning, udgivet af Jean Pio, H. Trier og P. Voss, 1879 (Kjøbenhavn: P. G. Philipsens Forlag), pp. 69-90.
%  Antiqua.  Quotations in German, French and English.
%
%  TRANSCRIPTION -- GENERATED; DO NOT EDIT BY HAND.  Høffding's three pieces of 1879-82 (Filosofien i Sverig; Om
%  Pessimismen; the review of Martensen's Jacob Böhme) are built by ONE pipeline (ebook-method/hoeffding-1879-82/) into a
%  combined working file, which split.py cuts into three transcriptions.
%  Rebuild from the repo root:   sh ebook-method/hoeffding-1879-82/rebuild.sh
%
%  STATUS: COMPLETE (2026-10-15): the volume's title page and the whole text, pp. 69-90, with its 5 notes.
%
%  EDITION.  The only printing.  Text of record: Nasjonalbiblioteket's scan 5913f0beb5001cb09055a914b1780933 (URN:NBN:no-nb_digitidsskrift_2020102381002_001,
%  ~400 ppi colour JPEG); extract ~/bibliotek/Høffding, Harald/1879-vor-ungdom-pp69-90.pdf (SCANS.tsv, sha256), made by extract.py
%  (title, Indhold, the article).  Page range from the scan: p. 69 opens the article, p. 90 ends it signed, p. 91 opens
%  „Om Rhythmens Betydning i Forhold til Opdragelsen“.
%
%  ACCURACY (2026-10-15; one collation for all three texts, 68 printed pages, ~28,100 words).
%  (a) No second copy of any of the three was found (Internet Archive, Nasjonalbiblioteket, Project Runeberg, Google Books;
%      HathiTrust not reachable).  The DRAFT is our own tesseract reading (tessdata_best, 300 dpi; dan+swe for Nordisk
%      tidskrift, dan for the others), chosen by measurement (words unattested in the corpus, and for Nordisk tidskrift
%      a majority of three readings: measure.py), laid over the working PDF and turned into LaTeX by draft.py.  The
%      WITNESS, collated word by word (collate/check.py): NB's ALTO (ABBYY) for Nordisk tidskrift and Vor Ungdom, the
%      Archive's tesseract for Ude og Hjemme.  1,213 disagreements in all.
%  (b) Rules (rules.py: ALTO reads æ as e, ø as g/e, the closing “ as *; glyph faults against attested words).
%      Audit of rule decisions: 0 contradicted.
%  (c) Every other disagreement read at the image BLIND (rounds b1/b2, the same readers and controls as Filosofien i
%      Sverig), plus the spots of the unattested-word and speck sweeps.  Italics read blind (i1).  OPEN 0.
%  (d) Every FOOTNOTE read whole at the image (reader N1), set from the reading; every note MARK (raised figures ¹))
%      placed as read and checked at a zoom.  STRUCTURE read at the image (reader K): title page, head, ending, signature.
%  (e) LETTERSPACING measured (spacemeasure.py) and judged at the image by reader K.
%  (f) QUOTATION MARKS: „...“ (every crop, the notes); both OCRs read „ as , and “ as % * ”: set by class (quotes.py).
%  PAGE CHECK.  SAMPLED PAGE CHECK (2026-10-15): 4 pages drawn at random from four strata (Filosofien i Sverig pp. 35, 60; Om Pessimismen
%      p. 83; Ude og Hjemme p. 366, left column), 1,595 words transcribed blind at the image by a reader who saw no other
%      text (pagecheck.py): 1 word error (ɔ: read as o:, then mended in all five places of the class), 1 punctuation error
%      („Følelse“. read Følelse,), 2 italic runs (a posteriori, a priori) -- all mended; the reader's 2 other differences were his
%      own misreadings at a zoom.  About 2 errors per 1,000 words before mending.
%  PRINTED AS IS (misprints kept, marked at the line):
%    p. 72: „Tyngdepuukt“ for Tyngdepunkt; blind readers (b1, two crops: "printed uu, clear")
%    p. 73: „menneskeskelige“ for menneskelige (menneske-/skelige); zoom z-marks2.png
%  CONVENTIONS.  Page turns: % --- p. N --- and \opage{N} (the printed page); a word divided over a page turn is
%  split with %.  Footnotes as \footnote at their marks, numbered per page as printed (¹)).  \emph{} =
%  letterspacing, \textit{} = italics.  The running heads and the sheet signatures are not transcribed.
% =====================================================================
\documentclass[12pt,a4paper,openany]{book}
\usepackage[utf8]{inputenc}
\usepackage[T1]{fontenc}
\usepackage{libertinus}
\usepackage{libertinust1math}
\usepackage[danish]{babel}
\usepackage{textalpha}   % Greek typed directly
\usepackage[protrusion=true,expansion=true]{microtype}
\usepackage{setspace}
\usepackage[margin=1.3in]{geometry}
\usepackage{fancyhdr}
\usepackage{hyperref}
\hypersetup{pdftitle={Om Pessimismen}, pdfauthor={Harald Høffding}, hidelinks}
\pagestyle{fancy}
\fancyhf{}
\fancyhead[LE,RO]{\thepage}
\fancyhead[C]{\textit{H. Høffding: Om Pessimismen (1879)}}
\fancypagestyle{plain}{\fancyhf{}\fancyhead[LE,RO]{\thepage}}
\setstretch{1.3}
\usepackage{marginnote}
\usepackage{graphicx}
\DeclareUnicodeCharacter{0254}{\reflectbox{c}}   % ɔ (id est)
\makeatletter
\renewcommand{\thefootnote}{\textsuperscript{\arabic{footnote}})}   % the print marks each page\'s notes with raised figures, ¹)
\makeatother
\newcommand{\opage}[1]{\marginnote{\footnotesize\textit{[#1]}}}
\newcommand{\signature}[1]{\par\smallskip\hfill #1\hspace*{2em}\par}
\pdfstringdefDisableCommands{\renewcommand{\footnote}[2][]{}\renewcommand{\opage}[1]{}}

\begin{document}

\begin{titlepage}
\begin{center}
{\Huge \emph{VOR UNGDOM}}\\[0.6em]
{\large TIDSSKRIFT}\\[0.2em]
{\scriptsize FOR}\\[0.2em]
{\large\bfseries\sffamily OPDRAGELSE OG UNDERVISNING}\\[0.8em]
{\scriptsize \emph{UDGIVET AF}}\\[0.3em]
{\emph{JEAN PIO, H. TRIER}}\\[0.2em]
{\scriptsize OG}\\[0.2em]
{\emph{P. VOSS}}\\[0.6em]
\rule{6em}{0.4pt}\\[0.3em]
{\emph{1879}}\\[0.2em]
\rule{6em}{0.4pt}\\
\vfill
% [an oval engraved portrait]
{\large\bfseries\sffamily KJØBENHAVN}\\[0.3em]
{\emph{P. G. PHILIPSENS FORLAG}}\\[0.3em]
1879
\end{center}
\end{titlepage}
\clearpage

% --- p. 69 ---
\setcounter{footnote}{0}%
\opage{69}\begin{center}{\Large\sffamily\bfseries \emph{Om Pessimismen.}}\end{center}\medskip

Det gaar med den aandelige Sundhed som med den
legemlige, at den lægger sig for Dagen ved, at man slet
ikke faar Anledning til at tænke paa den. Sundheden er
kun det nødvendige Grundlag for fuld Hengivelse til og
Opgaaen i Virksomhed og Arbejde. Og da Lykken nu
vist nok væsentlig bestaar i en saadan fuld og uhindret
Leven og Virken, er den altid for en stor Del ubevidst.
Livskraften udfolder sig, uden at man har nogen Grund
til at standse og reflektere over sin Tilstand eller sin Følelse.
Selv om der er Modstand og Hindringer, --- saa
længe de ikke staa i Misforhold til den Kraft, der er
paa rede Haand, ville de netop bidrage til at forøge og
forhøje Livsudfoldelsen. Den største Del af Menneskene
fører maaske sit Liv under saadanne Forhold. Der er
intet saadant dybt Brud i deres Natur, ingen saadan nagende
Tvivl, som kunde hindre Livsinstinktet og Selvopholdelsesdriften
fra at fortsætte deres Virksomhed. Deres
Virken, Nyden og Liden er bestemt ved en lille, overskuelig
Kreds af Forhold og Opgaver, og de have ingen
Grund til at fælde nogen Dom om Livet som Helhed.
Selv om enkelte store og dybe Sorger i nogle Tilfælde
kaste mørke Skygger ind over større eller mindre Livsafsnit,
% --- p. 70 ---
\setcounter{footnote}{0}%
\opage{70}saa formaa de dog ikke vedvarende at skade Livsfølelsen
i dens inderste Rod.

Paa forskjellige Maader kan den i selve Livsinstinktet
og Selvopholdelsesdriften liggende optimistiske Tendens
blive brudt. --- Ulykke og Lidelse kunne overvælde den
menneskelige Modstandskraft og gjøre fri, harmonisk Udvikling
umulig. Det enkelte Menneske, hvem dette vederfares,
vil føle Livet som Byrde og Smerte, og han vil
maaske udstrække sin Dom til det menneskelige Liv i
Almindelighed. Dog vil en dyb og fælles Følelse af Livet
som ulykkeligt ikke opstaa uden paa Grundlag af stærk
og rodfæstet Sympathi. Saa længe Ulykken kun rammer
den enkelte, eller de enkelte hver paa sin Maade, vil den
endnu staa som en Tilfældighed. Men naar Mennesket
tager andres Synd og Ulykke paa sig gjennem Sympathi
og Medlidenhed, begynder Ondet at gnave paa Livets
Rod. Derfor se vi ogsaa, at hvor Opfattelsen af Livet
som et Onde udbreder sig, har der samtidig fundet et
mægtigt Opsving af de sympathetiske Følelser Sted. Buddhismen
er ikke blot Pessimismens, men ogsaa den dybe,
inderlige og altomfattende Medlidenheds Religion. Kristendommen,
i Følge hvilken Verden ligger i det onde, er
tillige den almindelige Menneskekjærligheds Religion. I
den moderne pessimistiske Filosofi finde vi --- i hvert Tilfælde
i Theorien --- ligeledes de to Momenter forenede.
Schopenhauer lærte, at Verden var den sletteste af alle
mulige Verdener, og samtidig opstillede han Medlidenheden
som Moralprincip.

En pessimistisk Livsopfattelse forudsætter derfor stedse
en forholdsvis høj aandelig Udvikling. Uden Erindring,
uden udbredt Erfaring, uden Evne til at sætte sig paa
andre følende Væseners Stade kan den smertelige Følelse
ikke naa nogen høj Grad af Styrke og ikke komme til at
bestemme den Værdi, der tilskrives Livet. Det umiddel%
% --- p. 71 ---
\setcounter{footnote}{0}%
\opage{71}bare Livsinstinkt er Lidelsen for stærkt; det vil stedse
rejse sig paa ny, --- og selv om det ligger under i Praxis,
vil det hævde sig i Tanken: saaledes se vi ved Oldtidens
Slutning den antikke Optimisme samle sine sidste Kræfter
i Plotins Skrift mod de kristne Gnostikeres Pessimisme og
Dualisme.

Spørge vi nærmere, hvad det er for Ulykker, der saaledes
kunne gaa Menneskeheden til Marv og Ben, at Livsdriften
saares i sit inderste, da afgiver vist nok netop
\emph{Oldtidens Slutning} et oplysende Exempel. Naar en
stor og glimrende Kulturperiode er ved sin Fuldendelse,
naar alle Mulighederne synes at være opbrugte, naar der
ligesom er en Pavse i Verdenshistorien, da vaagner naturlig
Reflexionen, da føles Mangelen paa Livskraft. Den
gamle Pan er død, og man aner ikke, at hans Efterfølger
ligger i Vuggen. Ingen kunde tænke sig, at der i en
sekterisk Bevægelse hos et af Romerrigets mest ubetydelige
Folk og i de barbariske Nationer, der fra Tid til
anden løb Storm mod Rigets Grænser, laa Spirerne til
ny, frisk og storartet historisk Udvikling. Sociale Spørgsmaals
(saa vel som mange andre Spørgsmaals) Løsning ligger
ofte, hvor man mindst venter det.

Det gjælder --- som afdøde Sibbern saa ofte hævdede
--- om al Udvikling, at den foregaar sporadisk, begynder
ud fra spredte Punkter, og først naar de tilsyneladende
isolerede og tilfældige Udviklingsrækker skyde sammen og
forene sig i samme Løb, bliver det muligt at forudse Resultatet.
Mange Skuffelser opstaa, fordi man ensidig har
fulgt og hengivet sig til en enkelt Udviklingsrække og sat
alt ind paa den, skjønt den først i Forening med andre
Tendenser kan virke til blivende Fremgang. To af de
største pessimistiske Kriser i nyere Tid opstod paa denne
Maade, nemlig de, som betegnes ved Rousseaus og
Malthus's Optræden.

% --- p. 72 ---
\setcounter{footnote}{0}%
\opage{72}Den moderne Kultur har sit Tyngdepuukt i den intellektuelle % PRINTED AS IS (Tyngdepuukt: for Tyngdepunkt; blind readers (b1, two crops: „printed uu, clear“))
Udvikling. Det var paa Erkjendelsens Omraade,
at Udviklingens Lov først klart gik op for Bevidstheden.
Allerede Pascal sagde det berømte Ord: „La suite \textit{des
hommes, pendant le cours de tant de} siècles, \textit{doit être considérée
comme un} même \textit{homme qui subsiste} toujours
\textit{et qui apprend} continuellement“ („Menneskeslægten, hvis
Liv udfylder saa mange Aarhundreder, maa betragtes som
et og samme Menneske, der stedse er til og uafladelig
lærer noget nyt.“ \textit{Fragment d'un traité du vide}). Hos
dem, der i det attende Aarhundrede arbejdede for Fremskridtet,
blev derfor naturligen Oplysning Løsenet. Den
stigende Indsigt skulde omdanne Livet; til Dels havde
den begyndt at gjøre det. Den skulde befri fra Fordomme
og forældede Avtoriteter og bringe Naturen i Menneskets
Tjeneste. Over for dem, der mente, at alt var godt, naar
Folk blot lærte noget, stillede saa \emph{Rousseau} sit bekjendte
Paradox, at Kunst og Videnskab vare Menneskehedens
Ulykker, og sin Opfordring om at vende tilbage
til Naturtilstanden. Følelsen var for ham ét og alt. Han
har --- som Hettner siger --- befriet det menneskelige Gemyt
og indsat det hele og fuldstændige Menneske i dets ubrødelige
Ret, stillet det følende Hjærtes fremtrængende Sprog
mod det udvortes og snorrette i den ensidige Forstandsdannelse.
Men det var en skarp Dom, han fældede over
den tidligere Udvikling, og baade han og de Aander, hvis
Forgænger han blev, følte smertelig den Disharmoni, han
havde paapeget. Skjønt Erkjendelse og Følelse udspringe
af samme Rod, godtgjør Psychologien dog, at deres Udviklingsbetingelser
ikke ere de samme. Erkjendelsen gaar
hurtigere frem end Følelsen; Tanke forbinder sig lettere
med Tanke, end Følelse med Tanke; en Tanke gaar lettere
over til en anden Tanke, end en Følelse forvandles til en
anden. Følelsen kan altsaa ikke holde Skridt. Der gaar
% --- p. 73 ---
\setcounter{footnote}{0}%
\opage{73}lange Tider, fulde af Splid og Smerte især for dem, hos
hvem Følelseslivet har en særegen Inderlighed, før den ved
den raske intellektuelle Udvikling brudte Harmoni kan gjenoprettes.
Det var en stor og dyb Trang i den menneskeskelige % PRINTED AS IS (menneskeskelige: for menneskelige (menneske-/skelige); zoom z-marks2.png)
Natur, et Forhold af stor psychologisk og historisk
Betydning, Rousseau pegede hen paa ved sin pessimistiske
Dom over den moderne Kultur. Ogsaa her viste det naturlige
Slægtskab mellem Pessimismen og Sympathien sig.
De optimistiske Oplysningsmænd vare Aandsaristokrater;
deres Begreb om Dannelse kunde kun fyldestgjøres af en
lille Kreds, ud over hvilken de ikke rettede deres Blik;
Rousseau gik dybere og derfor ogsaa videre.\footnote[1]{\emph{Kant}, som i sine første Skrifter havde hyldet den intellektuelle Optimisme, blev stærkt paavirket af Rousseau's Skrifter. Han forkastede sin tidligere Anskuelse som Aandsaristokrati, og kasserede en Afhandling, han havde skrevet om Optimismen. Grunden blev herved lagt til hans senere ethiske Anskuelse. (Se \emph{Dieterich}: Kant und Rousseau. Tübingen 1878. p. 7---13; 97---98).} Men
i sin Iver opfattede Rousseau Tænkningen, Reflexionen,
som Frugt af et Syndefald. „L'état de réflexion est un
état contre la nature, et l'homme qui médite, est un
animal dépravé“ („At tænke strider mod Naturen, og
det Menneske, der grubler, er et vanslægtet Væsen.“
Discours \textit{sur l'origine de l'inégalité parmi les hommes).}
Det er sikkert, at Tanken fører ud over den Tilstand,
hvor Livsinstinktet raader umiddelbart, hvor ingen Muligheder,
ingen Sammenligning, ingen Idealer gribe forstyrrende
ind. Men Rousseau, denne Følelsens Herold, modsiger
sig selv, naar han i Modsætning til Reflexionen anbefaler
sin saa kaldte Naturtilstand, hvor „Indbildningskraften intet
maler og Hjærtet intet fordrer“. Sagen er den, at skjønt
den intellektuelle Udvikling foreløbig volder Disharmoni
og gjør Brud paa det instinktmæssige Følelsesliv, er den
% --- p. 74 ---
\setcounter{footnote}{0}%
\opage{74}dog selv en Betingelse for Følelseslivets højere Udvikling.
Følelsens Højde, Renhed og Omfang beror paa, hvor mange
Tankeelementer den har optaget i sig. Haab og Frygt,
Kjærlighed og Had, Beundring og Enthusiasme opstaa kun
ved Sammensmeltning af Erkjendelses- og Følelseselementer.
Ved at fordømme Tankelivet berøver man altsaa tillige en
hel Række af de højeste og ædleste Følelser Betingelserne
for deres Udvikling. Rousseau opfatter kun Erkjendelsen
fra den Side, hvor den træder i Modsætning til den umiddelbare
Følelse og, for en Tid i det mindste, bringer
Splid og Uro i Stedet for den trygge Hvilen i Instinktet.
Naar han forlanger Tilbagegang til Naturtilstanden, fordrer
han egentlig, ligesom de moderne Pessimister, Tilbagegang
til det ubevidste. Forudsætningen er: Lykken er ubevidst;
Selvbevidstheden gjør os ulykkelige.

Ved Siden af Oplysningen var der et andet Middel,
hvorpaa det attende Aarhundrede byggede sine Forhaabninger
om ubegrænset Fremskridt og Lykke for Menneskeheden:
nemlig politiske Reformer. Naar blot de ydre
Hindringer for almindelig Lykke, der laa i Samfundets
Former og Institutioner, bleve fjærnede, vilde alt blive
godt. I sin mest overstadige Form blev denne Optimisme
udtalt af Condorcet og Godwin; men den er herskende i
Litteraturen hen imod Aarhundredets Slutning (smlgn. Greg:
\textit{Enigma's of Life.} 4 ed. \textit{London 1873.} p. 54). Da
kom som et Lyn fra en klar Himmel Malthus's berømte
Bog om Befolkningsspørgsmaalet, i hvilken han viste, at
Menneskeslægten tiltager stærkere end Ernæringsmidlerne,
saa at de sociale Ulykkers Kilde ikke maa søges paa
Overfladen, men i selve Grundlaget for det menneskelige
Liv. Ogsaa her var den pessimistiske Anskuelse den dybest
gaaende. Siden den Tid har den sociale Pessimisme adskillige
Gange kæmpet mod den politiske Optimisme, og
denne Disharmoni vil --- ligesom hin beslægtede psycho%
% --- p. 75 ---
\setcounter{footnote}{0}%
\opage{75}logiske Disharmoni --- kræve lange Tider til sin Løsning. I
begge Henseender er det af største Betydning, at Følelsen
er vakt for, at der er en Kamp, som maa bestaas. --- Den
sociale Pessimisme synes, ligesom den psychologiske, at
kræve Tilbagevenden til det ubevidste. De socialistiske
Systemer have en mere eller mindre reaktionær Karakter,
stille sig i et mere eller mindre polemisk Forhold til den
individuelle Frihed, der paa en Maade er i Samfundet hvad
Selvbevidstheden, Intelligensen, er i det enkelte Menneske.
Og dog er det sikkert, at der i den individuelle Frihed,
som saa lange Tiders Kampe har banet Vejen for, er
naaet et virkeligt Gode, hvilket er mere, end man kan sige
om de hidtil udkastede sociale Fremtidsidealer. At begynde
med at kue den er lige saa misligt som at ville
frelse Følelsen ved at binde Forstanden.

Det er ikke umuligt, at den stærke pessimistiske
Strømning i \emph{Tyskland i vore Dage} for en stor Del
skyldes denne Modsætning mellem det politiske og det sociale.
Efter at det store politiske Maal var naaet, som
havde været Gjenstand for saa mange Forhaabninger og
Anstrængelser, og som havde kostet saa mange Lidelser,
maatte de sociale Misforhold træde saa meget des skarpere
frem. Lyset gjør Mørket mere sort. Med Opnaaelsen af
et stort Maal er der tillige hyppig forbundet et vist Vemod.
Man føler Tomhed og spørger: hvad nu? Det, der glimrede
som det fjærne Maal, mister let Glansen, naar det
gaar over til Fortiden. Det hører Følelseslivet til at
bevæge sig gjennem Modsætninger; tages Skranke og Modstand
bort, føler Mennesket sig som i det tomme Rum.
\emph{Stuart Mill} havde i sin Ungdom en Krise at gjennemgaa,
som i denne Sammenhæng er af Interesse. Den Livsanskuelse,
i hvilken han tidlig var bleven ført ind, havde
lært ham „at grunde sin Lykke paa et varigt og fjærnt
Maal, hen imod hvilket der bestandig kunde gjøres Frem%
% --- p. 76 ---
\setcounter{footnote}{0}%
\opage{76}skridt, uden at det dog nogen Sinde kunde udtømmes ved
fuldstændig Opnaaelse“. Men i sin Begejstring for Fremskridt
og Udvikling standsedes han i en mørk Time ved
det Spørgsmaal: „Sæt, at dit Livsmaal blev opfyldt; sæt,
at alle de Forandringer i Institutioner og Anskuelser, som
du attraar, nu i dette Øjeblik bleve fuldstændig virkeliggjorte;
--- vilde du saa være inderlig glad og lykkelig?“
Og han maatte besvare dette Spørgsmaal med Nej. „Modet
svigtede mig,“ siger han; „hele den Grundvold, hvorpaa
mit Liv var bygget, styrtede sammen.“ For enhver Livsanskuelse
gives der tilsvarende Anfægtelser, hvor selv det,
der ellers staar som det højeste og herligste, mister sin
Værdi og alt bliver tomt og øde. Stuart Mill indsaa, at
han hidtil havde indrømmet Følelsen for ringe en Rolle;
men det var dog ikke blot det, der førte ham ud over
Krisen. Han gjennemtrængtes af den Overbevisning, at
Mennesket maa underordne sig et Maal, en Opgave, over
for hvilken det kan glemme sig selv. Betragter man sig
selv som Formaal, vil man føle Tomheden. „Gjør dig
det Spørgsmaal, om du er lykkelig, --- og du vil ophøre
at være det. Man maa ikke betragte Lykken, men et
Formaal, der ligger ud over den, som Livets Opgave. Lad
din Selvbevidsthed, din Efterforsken, din Selviagttagelse
gaa ud herpaa, --- og du vil\dots{} indaande Lykke med
den Luft, der omgiver dig“ (John Stuart Mill's Levned.
Kbhvn. 1874. p. 136; 145).\footnote[1]{Stuart Mill nærmede sig her, som han selv bemærker, \emph{Carlyle}'s „Anti-Selvbevidstheds Theori“. Thi ogsaa denne originale Aands Vidnesbyrd kan anføres for, at fuld Bevidsthed --- især den fulde Bevidsthed om Opnaaelsen af et Maal --- er Lykkens Fjende. „Selvbetragtning“, siger han, „er stedse et Tegn paa Sygdom og Egoisme. Det er usundt at fortære sig selv i tom Anger og Frygt, og endnu usundere at svælge i Betragtningen af sin egen Handlen. Det er til ingen Nytte at maale den Vej, man har tilbage%
\opage{77}lagt; det kommer an paa at gaa frem. Er den fornuftige Vilje, der bor i vort allerhelligste, virkelig fri som en Guddom, saa vil den fuldkomne Lydighed være tavs. --- --- Ubevidsthed og Hemmelighed ere Tegn paa alt stort; alt, hvad man fuldstændig kan gjennemskue, er trivielt. Med Rette gjorde de gamle Tavsheden til en Gud; med Rette synge nyere Digtere Hymner til Natten.“ (\emph{Julian Schmidt}: Portraits aus dem 19ten Jahrhundert. 1878. p. 107 f.)} Opnaaelsen af et Formaal
% --- p. 77 ---
\setcounter{footnote}{0}%
\opage{77}kan gjøre ulykkelig ved at rette Reflexionen mod det egne
Selv, der hidtil var traadt tilbage for Opgaven. I Bestræbelsen,
den fremskridende Stræben, maa Lykken søges.
Derfor kunde Leibnitz ikke tænke sig Saligheden anderledes
end som bestandig Fremskriden; ellers vilde den medføre
Sløvhed \textit{(nisi beatitudo in progressu consisteret,
stuperent beati. Epp.} ad Wolf).

Den tyske Pessimisme kan dog næppe fuldstændig
forklares af den Krise, et stort Maals Opnaaelse fører med.
Der er sikkert mange andre Aarsager, som have medvirket
til, at det er kommet saa vidt, at Bismarck i Fjor kunde
sige i den tyske Rigsdag: „De tyske have stærk Tilbøjelighed
til Utilfredshed; jeg véd ikke, om nogen af os har
set en tilfreds Landsmand?“ Dr. Waldstein, som særlig har
undersøgt Spørgsmaalet om den tyske Pessimismes Aarsager,
og af hvem en Artikkel har været trykt i „Nær og
Fjern“ for 1878, mener, at baade Undervisningsvæsenets
og det selskabelige og sociale Livs Mangler og Ensidigheder
have bidraget meget til dette Tidsfænomen. Mange
af Aarsagerne ere fælles for alle Lande, og ville ved
nærmere Undersøgelse kunne føres tilbage til de Vendepunkter,
der betegnes ved Rousseau og Malthus. Vi ere
endnu midt i de Kriser, disse Navne pege hen paa. Vi
mangle Harmoni i vor Livsanskuelse, Tillid og Tryghed
med Hensyn til det offentlige Livs sociale Grundlag.

Man maa ikke glemme, at Pessimismens 
Hovedord%
% --- p. 78 ---
\setcounter{footnote}{0}%
\opage{78}fører, \emph{Arthur Schopenhauer}, ikke er nogen moderne
Forfatter i Ordets snævrere Betydning. Hans Ungdom
falder i Romantikkens Periode, hans Hovedværks første
Del udkom 1819, og han var da 31 Aar gammel. Det
er fra flere Sider blevet fremhævet, at den pessimistiske
Poesi (Byron, Leopardi, Heine) gik forud for den pessimistiske
Filosofi (Schopenhauer, Hartmann). Den pessimistiske
Poesi hænger nøje sammen med den fra Rousseau
udgaaende Følelsesretning. Pessimismens engelske Kritiker,
James Sully, har i en særegen Afhandling gjort opmærksom
paa, at Digterne til alle Tider, ogsaa i Oldtiden,
dvæle ved Livets sørgelige Side. Han søger Aarsagen
hertil i, at det er vigtigere for de levende Væsener at
udtrykke deres Smerte end deres Lyst, hvorfor Instinktet
til at lade Smerten faa Udbrud naturlig vil være det
stærkest udviklede. Nu kan vel det civiliserede Liv siges
at have en Tendens til at vende Forholdet om; i godt
Selskab maa man jo ikke græde eller jamre sig; og
den moralske Styrke lægger sig for Dagen i, at man kan
dæmpe Sorgen. Men den oprindelige Drift lader sig ikke
kue, især ikke hos Digteren, der har en finere, sensiblere
Konstitution end andre Mennesker, og hos hvem Følelse
og Lidenskab spille saa stor en Rolle. Der ligger en stor
Trøst og Lettelse i at give Lidelsen et stærkt og heftigt
Udtryk. Desuden næres Indbildningskraften ved den virkelige
Verdens Ufuldkommenheder; og omvendt kommer
Verden netop til at staa som mørk og ufuldkommen over
for Fantasiens Idealer \textit{(Revue philosophique.} Avril 1878).

Der kan ikke drages nogen skarp Grænse mellem den
pessimistiske Poesi og den pessimistiske Filosofi. Hvis
man ved Filosofi forstaar den kritisk-videnskabelige Undersøgelse
af vor Tænknings Grundbegreber, er Filosofien
hverken optimistisk eller pessimistisk. Dens Resultater
kunne til en given Tid synes at pege i den ene eller den
% --- p. 79 ---
\setcounter{footnote}{0}%
\opage{79}anden Retning; men dette er i sig selv uvæsentligt for
den. Anderledes naar man tager Ordet Filosofi i en videre
Betydning. Enhver Livsanskuelse, der med klar Reflexion
søger at føre Livets indviklede Forhold tilbage til
visse Principper, kan, især naar disse Principper igjen
søges bragte i Forbindelse med hin kritisk-videnskabelige
Undersøgelses Resultater, kaldes Filosofi. I de Systemer,
Filosofiens Historie fremstiller for os, finde vi gjærne begge
Bestræbelser nøje forbundne, ofte indfiltrede i hinanden.
Filosofiens Stræben gaar i vor Tid ud paa at gjøre den
kritisk-videnskabelige Undersøgelse til Hovedsagen for
derved at vinde et fast Underlag og Udgangspunkter for
sikker Fremgang. Schopenhauer indtager i denne Henseende
en ejendommelig Stilling. En stor Del af hans
Undersøgelser tilhører afgjort den videnskabelige Filosofi, og
for en ikke ringe Grad har dette været Grunden til hans
Anseelse under Reaktionen mod den spekulative Filosofi.
Han søgte at arbejde videre paa det af Kant givne erkjendelsestheoretiske
Grundlag. Hans Filosoferen har en
analytisk Karakter: fra det i Erfaringen givne søger han
tilbage til Elementerne, til de Urfænomener, der betegne
Erkjendelsens Grænser. Her har han et fast og sikkert videnskabeligt
Stade over for den spekulative Filosofi (Fichte,
Schelling, Hegel), der vilde udlede alt ved apriorisk Tænkning
af ét eneste absolut Princip. Men denne Besindighed
holder hos ham ikke ud til Enden. Paa et vist Punkt
aabner han Gærdet, opgiver Ævret og lader den af Følelsen
inspirerede Fantasi raade. „Tingen i sig selv“,
som Kant havde erklæret for det evig ubegribelige, skulle
vi i Følge Schopenhauer umiddelbart\footnote[1]{Dog er Schopenhauer her noget vaklende; undertiden synes han at betragte Bevidstheden om Viljen som middelbar. Han siger f. Ex. Parerga und Paralipomena. Berlin 1851. I p. 18: „Weil doch wir selbst unstreitig real sind, muss \opage{80}aus dem Innern unsers eigenen Wesens die Erkenntniss des Realen irgendwie zu schöpfen sein. In der That nun tritt es hier, \emph{auf eine unmittelbare Weise}, in's Bewusstsein, nämlich als Wille.“ Men II p. 79 hedder det: „\emph{Wäre dieses Sichbewusstwerden} [des Seienden selbst als Wille] \emph{ein unmittelbares}, so hätten wir eine völlig adäquate Erkenntniss des Dinges an sich. Weil es aber dadurch \emph{vermittelt} ist, dass etc.“} blive os bevidste,
% --- p. 80 ---
\setcounter{footnote}{0}%
\opage{80}i det vi ved at vende Blikket ind ad mod os selv ligesom
mærke Tilværelsens Puls slaa i os, føle Urkraften i alt
røre sig som en ubændig Drift, en Higen og Villen, der
aldrig fuldt tilfredsstilles. Denne i sig selv ubevidste
Vilje er for Schopenhauer Princippet i Tilværelsen; alle
Naturkræfter og Naturformer ere dens Ytringer, betegne
Skridt paa dens Vej. Bevidstheden og Tanken er kun,
opstaaet som et Middel for Viljen til at opnaa sit Formaal
tilfredsstille sin Hunger og Tørst efter Existens. Men
med Bevidstheden opstaar ogsaa Smerten. Den ubevidste
Vilje er en Brynde, der kun foreløbig mættes, men altid
melder sig paa ny. Al Lyst er kun negativ, er en Illusion,
der opstaar, naar den hede Attraa et Øjeblik slukkes;
Smerten er derimod positiv, det egentlige Udtryk for Urviljen,
naar denne bliver Gjenstand for Bevidsthed. Det
gjælder derfor fremfor alt at ophæve Viljen til Livet for
at stoppe Lidelsens Kilde. Dertil staa forskjellige Veje
aabne: kunstnerisk Anskuelse og videnskabelig Tænkning,
hvor Mennesket glemmer sig selv i den rene Kontemplation
uden nogen Attraa eller Villen, --- Medlidenheden, i
hvilken Individet gjør sig ganske til ét med andre Individer
og saaledes ophæver sin Vilje, --- Resignation og Askese,
der paa sit Højdepunkt endog kan føre til, at Mennesket
ikke længer tilfredsstiller det vegetative Livs Fordringer.
Viljen behøver altsaa „Kvietiver“, ikke Motiver; den skal
dysses hen som et sygt og uroligt Barn, ikke plejes og
udvikles.

% --- p. 81 ---
\setcounter{footnote}{0}%
\opage{81}Denne Livsanskuelse som Schopenhauer fremstiller
med stor Energi og digterisk Evne, taaler næppe nogen
nærmere logisk Undersøgelse. Hos hans berømte Antipoder
(som han aldrig bliver træt af at overdænge med
Haansord), hos Fichte, Schelling og Hegel, maa vi netop
beundre den Kraft og Konsekvens, hvormed en ideal
Livsanskuelses Indhold er formet til et sammenhængende
System, med en bestemt Bestræbelse for at knytte Led til
Led. Men her hvor deres Styrke er, er netop Schopenhauers
Svaghed. Sagen er: den dybe pessimistiske Følelse,
som i Litteraturens Historie saa ofte slog ud i
lidenskabelige Udbrud, som især Rousseau og Digterne i
Aarhundredets første Del udtalte, uden at den dog altid
var deres sidste Ord, --- den er bleven ét og alt hos
Schopenhauer, har hos ham for første Gang i Tankens
Historie søgt at bemægtige sig Enevælden. Den gjør det
ved et Statskup. Thi der kan næppe paavises nogen nødvendig
Forbindelse mellem den kritisk-videnskabelige Del
af hans Filosoferen og hans Lære om Urviljen. Med de
samme almenfilosofiske Forudsætninger kunde Verdensgaaden
lige saa godt være bleven gættet paa en anden Maade.
Det er Følelsen, der har dikteret Løsningen.

Men ogsaa afset herfra har denne Løsning en stor
Interesse. Schopenhauer's Livsanskuelse træder i en afgjort
Modsætning til hele den evropæiske Menneskeheds
Livsopfattelse i dens forskjellige Former. Saa dyb en Følelse
af Livets Smerte og Synd der end kan være udtalt
af Religionsstiftere og filosofiske Tænkere, saa ligger den
stille Tro dog stedse til Grund, som en yngre Digter hos
os har udtalt i de Ord, at „oven over Skyerne er Himlen
altid blaa“. At Verdens Substans var god, derom
nærede man ingen Tvivl; derom er selv Augustinus med
sin Tro paa, at Flertallet af Menneskene er bestemt til
evig Fordømmelse, enig med Platon og Leibnitz. Og selv
% --- p. 82 ---
\setcounter{footnote}{0}%
\opage{82}hvor man havde løsrevet sig fra den platoniske og den
kristelige Tradition, troede man dog paa, at Livets Kjærne,
om den end var nok saa bitter, kunde udfolde sig og blive
til Frugt. Men for Schopenhauer strækker Raaddenskaben
sig helt ind i Verdens inderste Væsen; Substansen er
slet, og al Udvikling er en Illusion. Optimismen er for
ham en Ryggesløshed, som især Jøderne og de moderne
Demagoger have udbredt. Det eneste, der forsonede ham
med det gamle Testamente, var Fortællingen om Syndefaldet,
og som Kristendommens egentlige Kjærne betragtede
han Læren om Arvesynden og Fordømmelsen. Derimod
optog han med stor Sympathi og Begejstring den buddhistiske
Lære om Tilværelsens Smerte og om Nødvendigheden
af at vende tilbage til Nirvana, det uendelige Intet. Han
føler sig i bestemt Strid med den evropæiske Livsanskuelse,
der baade i sine religiøse og sine humane Former opfatter
Livet som Udvikling og Fremskridt og stiller Friheden
over Nydelsen og Sikkerheden. „Asiaterne“, siger Draper
(Den evropæiske Aandsudviklings Historie. Kbhvn. 1875.
p. 52), „have vejet Frihed og Sikkerhed mod hinanden og
med Overlæg foretrukket denne og overladt det til Evropa
at sukke efter hin. Friheden skattes udelukkende under
et virksomt Liv, men Asiens Liv er væsentlig af passiv
Natur, og det higer efter Ro.“ Kvietiver, ikke Motiver,
er det jo, Schopenhauer søger.

Det kunde synes at tyde paa Energi i Livsanskuelsen,
naar Schopenhauer opfatter Verdens inderste Grund som
Vilje. Men han er netop her overordentlig vilkaarlig og
overfladisk. Medens han i det enkelte, især i sine mere
digteriske og populære Udviklinger, ofte røber fin psychologisk
Sans, mangler han i højeste Grad Præcision og
Grundighed ved Bestemmelsen af psychologiske Grundbegreber.
--- Han anser det for uvæsentligt ved Viljen, at
den ledes af bevidste Motiver; dette er kun Tilfældet hos
% --- p. 83 ---
\setcounter{footnote}{0}%
\opage{83}Mennesker og Dyr. Derimod skal det være den ubevidste
Vilje, der udøver alt i Naturen. Begrebet Kraft skal
underordnes Begrebet Vilje; thi hvad Kraft er, forstaa vi
kun ved Hjælp af Viljen, som vi umiddelbart ere os bevidste;
for at forstaa det ubekjendte maa vi føre det tilbage
til det bekjendte. Men hvorledes lægger Viljen sig
da for Dagen i vort Indre, saaledes at vi kunne blive os
den umiddelbart bevidst? Til Viljen henfører Schopenhauer
ikke blot alle Beslutninger, men ogsaa alle Drifter,
Ønsker, Lidenskaber og „alle de Bevægelser i vort Indre,
som man plejer at henføre under det vide Navn Følelse“.
Han skjelner altsaa ikke mellem Følelse og Vilje. I sig
selv er dette nu ogsaa et meget vanskeligt Punkt i Psychologien;
paa en enkelt Undtagelse nær (Alexander Bain)
findes der ingen Forarbejder til at komme til Klarhed over
dette Punkt. Men den Flothed, hvormed Schopenhauer
gaar til Værks, er exempelløs.

Naar man gjør Følelse og Vilje til ét, vil det kunne
komme af, at man reducerer Følelsen til Vilje, eller omvendt.
Schopenhauer, for hvem Viljen er ét med Brynde,
Higen, Drift, opfatter den kun som passiv, gjør altsaa i
Virkeligheden Viljen til Følelse. Dette er karakteristisk,
da vi have set, hvor væsentlig en Rolle Følelsen spiller,
hver Gang Pessimismen optræder. Mennesket staar hos
Schopenhauer overvejende passivt over for Yderverdenen;
 han søger kun at dække sig saa godt som muligt og ønsker
blot at faa Ende paa Legen saa snart som muligt. Hans
Kæmpen er i det højeste defensiv, aldrig offensiv. „Den
kaade Skæbnes Pileskud og Stenkast“\footnote[1]{I den berømte Monolog, i hvilken disse Ord forekomme, skildrer Hamlet, som \emph{Bucknill} (\textit{The mad folk of Skakspeare.} p. 86) bemærker, sin egen Karakter: „saa ivrig til at tænke, saa sløv til at handle, saa skarp i sin Vurdering af Livets Onder, saa utilbøjelig til at træde virksomt op imod dem.“} ramme ham, og
% --- p. 84 ---
\setcounter{footnote}{0}%
\opage{84}han taaler det med mere eller mindre Resignation. Viljens
Aktivitet maa falde bort, naar den løsrives fra sin
Forbindelse med Tanken om positive Formaal, der skulle
naas. Kun ved denne Forbindelse bliver Viljen mere end
uklar Brynde. Fichte, hvis utaknemmelige Discipel Schopenhauer
var, gjorde allerede Viljen til det absolute, men
det var netop som den uendelige Stræben efter Frihed,
den positive Virken for Virkeliggjørelse af Idealet.

Men ogsaa Følelsen lider Uret ved hin psychologiske
Uklarhed. Schopenhauer betragter Smerten som det egentlig
positive, medens Lystfølelsen stedse er negativ, i det
den opstaar ved, at Driften for en Stund slukkes. Hvorledes
der hos Schopenhauer i det hele kan blive Tale
om Lyst, naar Viljen er Danaidernes Kar, der aldrig fyldes,
er ikke godt at indse. Men i hvert Tilfælde fæster han
kun sit Blik paa de Lystfølelser, der umiddelbart ere knyttede
til Selvopholdelsen (f. Ex. Sultens og Tørstens Tilfredsstillelse).
Den blinde Drift klamrer sig til Livet, og
naar det lykkes, opstaar den Illusion, man kalder Lyst,
Glæde eller Lykke. Vi glæde os altsaa kun, naar vi have
ladet os narre af den ubevidste Vilje. I sig selv er der
ingen Realitet i Glæden eller Lykken. Livets største Goder,
Sundhed, Ungdom og Frihed, skatte vi først, naar vi have
mistet dem.

Det er sandt, at der i Driften hyppig er et Moment
af Ulyst. Men deraf følger ikke, at hele Lystfølelsen ved
Driftens Tilfredsstillelse beror paa, at denne Ulyst ophører.
Der er i Driften, saa længe den ikke afgjort er bleven
skuffet, tillige en Forudgriben, en Lyst i Haabet, og Lystfølelsen
ved Driftens Tilfredsstillelse beror til Dels paa,
at denne forud grebne Lyst udfolder sig til fuld Virkelighed
og Styrke. Selv ved de til Individets og Slægtens
Opholdelse knyttede Drifter sker der, naar de tilfredsstilles,
mere, end naar man lægger + 3 til --- 3. I Følge Scho%
% --- p. 85 ---
\setcounter{footnote}{0}%
\opage{85}penhauer gaar det lige op, og det er det resulterende 0,
der for os ved en Illusion staar som positiv Lyst. Til den
normale Udførelse af de organiske Funktioner er der
knyttet en lykkelig Grundstemning, en Følelse af Lethed
og Frihed, som vi ganske vist sjælden direkte henvende
vor Opmærksomhed paa og i Reglen kun mærke, naar den
afløser eller afløses af Ildebefindende (ligesom vi i Følge
Helmholtz opfatte Harmonien ved dens Modsætning til
Disharmonien). Men skulde det være saa stort et Gode,
om vi med udtrykkelig Bevidsthed nød vor Sundhed, Ungdom
og Frihed? Det gode ved disse Ting bestaar dog
netop i, at de sætte os i Stand til fuld Udvikling og Anvendelse
af vore Kræfter. En lykkelig Grundstemning i
Livsfølelsen er Betingelse for al vor Virksomhed; sygelig
Livsfølelse vil snart lamme alle vore Evner. Livsfølelsen
er Roden til hele vort aandelige Liv; Sindssygdom begynder
i Regelen med Forstyrrelse og Sygelighed i den.\footnote[1]{En tysk Læge har søgt at vise, at Schopenhauer led af en arvelig Sindssygelighed. Hans Bedstemoder blev sindssyg i sin Alderdom. Et af hendes Børn var Idiot. Schopenhauers Fader havde et mørkt og tungsindigt Temperament og led af voldsomme Vredesudbrud. Mod Slutningen af sit Liv viste han en saadan Ængstelighed med Hensyn til sine Forretninger, at man betragtede ham som sindssyg: hans Død ved at drukne i en Kanal blev af mange betragtet som Selvmord. --- Denne Ængstelighed og ubestemte Frygt var ogsaa et Hovedtræk i Arthur Schopenhauers Temperament. --- (Smlgn. J. \emph{Sully}: \textit{Pessimism. A history and a criticism. London 1877.} p. 75. 82.)}
Dens Betydning og positive Værdi beror ikke paa, at den
selv trænger sig frem i Forgrunden, men paa, hvad den
bærer og betinger.

Schopenhauer oversér dernæst ogsaa, at Lystfølelsen
ikke behøver at have Smerten til Baggrund, selv om den
træder skarpest frem i dette Modsætningsforhold. Den kan
% --- p. 86 ---
\setcounter{footnote}{0}%
\opage{86}have en svagere Lystfølelse til Baggrund. Det vilde sikkert
være den lykkeligste Existens, hvor forud gaaende Glæder
kunde tjene de følgende til Folie. Skranker for vor Virksomhed
fremkomme jo ikke blot ved, at en ydre Magt
griber ind i vor Existens og volder positiv Smerte, men
ogsaa ved, at der endnu er noget, som ikke er opnaaet,
skjønt det kan opnaas. Disse Skranker kunne lægge saa
betydeligt Beslag paa os, at der, uden egentlig forudgaaende
Ulyst, ved Arbejdets Gjennemførelse opstaar en positiv
Lyst. I selve Bestræbelsen er, saa længe den gaar frem,
--- saa længe Modstanden ej overskrider en vis Grænse, ---
Lystmomentet det overvejende. Men dette er et Ideal, der
kun delvis falder i vor Lod. Vore Livsvilkaar ere saadanne,
at Smerten idelig griber ind med sine mørke Skygger.
Og netop hos dybere udviklede Individer vil dette være
Tilfældet, da de ej blot bære deres egne Lidelser, men i
Sympathien ogsaa andres. Dog er det sikkert, at netop de
ædleste Nydelser ere hævede over Selvopholdelsesdriftens Uro
og Smerte. Schopenhauer maa indrømme dette om den Glæde,
Kunsten og Tænkningen volde; men han kommer derved i
Modsigelse med sig selv: thi i det rene Nirvana --- og
kun dér er „Viljen til Livet“ fuldstændig ophævet --- kan
der hverken gives Kunstanskuelse eller Tænkning.

Men Schopenhauer kan i Grunden ikke forklare Forskjellen
mellem ideelle og fysiske Nydelser. Som han ikke
anerkjender nogen historisk Udvikling i Slægten, saaledes
heller ikke nogen psychologisk Udvikling i Individet. At
der hos Mennesket udvikler sig et Tankeliv, det skal efter
hans Mening kun gjøre, at Mennesket rystes langt stærkere
af Lidenskaberne end Dyret; men det skal ikke øve nogen
væsentlig Indflydelse paa selve Drifterne og Følelserne:
„\emph{am Ende und im Realen} handelt es sich nur um die
\emph{selben Dinge}, die auch das Thier erlangt, und zwar
mit unvergleichlich geringerem Aufwande von Affekten
% --- p. 87 ---
\setcounter{footnote}{0}%
\opage{87}und Quaalen“ (Parerga und Paralipomena. Berlin 1851.
II p. 250 f). Dette er konsekvent, naar Viljen som blind
umættelig Drift rører sig i alt. Med stigende Bevidsthed
maa alt kun blive værre. Men det strider aldeles mod
den psychologiske Iagttagelse, der netop (som allerede
Platon og Stoikerne have vist) godtgjør Tankens lutrende,
udviklende og harmoniserende Indflydelse paa Følelseslivet.
Under det inderlige Vexelforhold, Psychologien paaviser
mellem Erkjendelse og Følelse, laaner hin sin Styrke fra
denne, men giver den til Gjengjæld Indhold, Omfang og
Klarhed. Som vi oven for have sét, koster det lange
Tider, inden Erkjendelsens Udvikling fuldt kommer Følelsen
til gode; men at det kan ske, og at den naturlige
Udvikling fører i den Retning, hører der stor Forblindelse
til at nægte. Det er muligt, at de brutale Selvopholdelsesinstinkter
(Schiller's Hunger und Liebe) endnu
væsentlig beherske Livet; men det ér i hvert Tilfælde et
betydningsfuldt Fremskridt, at de nu nødes til at godtgjøre
deres Berettigelse ved at knyttes til et mere eller mindre
ideelt og universelt Formaal (i Familien, Staten o. s. v.).
Det ideelle, der fra først af er Middel, bliver efterhaanden
Formaal.

Den aprioriske Antagelse af blind Drift og Brynde
som Livets Kjærne, en Antagelse, som selve den pessimistiske
Følelse dikterede, gjør det umuligt for Schopenhauer
at være en uhildet Iagttager af Virkeligheden. Man skulde
synes, at Erfaringen var den eneste Vej, ad hvilken man
kunde faa at vide, om der findes Lykke i Verden. Ud af
sit eget Indre kan man ikke hente Afgjørelsen af dette
Spørgsmaal; thi fordi jeg er ulykkelig, er det jo ikke sagt,
at alle andre ere det. Schopenhauer er snart færdig med
Erfaringen. „Dersom man“, siger han, „i Korthed vil
prøve den Paastand, at Nydelsen overvejer Lidelsen her
i Verden, eller at de i det mindste opveje hinanden, saa
% --- p. 88 ---
\setcounter{footnote}{0}%
\opage{88}sammenligne man det Dyrs Fornemmelse, der fortærer et
andet, med dette andet Dyrs Fornemmelse.“ Der er
naturligvis større Smerte ved at blive spist, end der er
Glæde ved at spise. Men saa simpelt er Regnestykket
ikke. Menneskene have jo i det mindste sikret sig nogenlunde
mod at blive spist; deres Liv behøver ikke at være
en stadig Frygt. Men Erfaringens Vidnesbyrd kan Schopenhauer
ikke tilskrive stor Vægt. Selv om man vilde forsikre
ham, at man var lykkelig, vilde han nægte, at det
var sandt. Han véd jo, hvad der ligger bag ved, --- at
det er „Viljen til Livet“, vi blindt følge. Men det er en
Anmasselse at kalde al Lyst for en Illusion. Om der
virkelig føles Lyst eller Smerte, kan kun Individet afgjøre,
og han lader sig ikke afdisputere sin Følelse. Derfor
hjælper f. Ex. Ræsonnement saa lidt over for Hypochondri, „Fornuften“ saa lidt over for Lidenskaben. Selv
den saa kaldte illusoriske, usande eller indbildte Glæde
er en virkelig Glæde. Hvad man mener med disse Udtryk
er egentlig det, at Følelsens Aarsag er en Indbildning
eller ikke har nogen Bestaaen. Kun ved at kritisere
Følelsens Aarsag og Indhold kan man komme den
selv til Livs.

De ældre Pessimister opfattede endnu det ubevidste
som lykkeligt. Den Lykke, de ikke fandt i det nærværende,
lagde de tilbage i en Guldalder, et Paradis eller
en barnlig Naturtilstand. Men den moderne Pessimisme,
som bilder sig ind at have opdaget Tilværelsens inderste,
opfatter det ubevidste Liv som Smerte og Uro. Den indvikler
sig derved i en af sine mange Modsigelser; thi paa
den anden Side skal jo netop Tilbagegang til det ubevidste
være Lægedommen for al Smerte.

Schopenhauer's Efterfølger, Eduard v. Hartmann,
har i flere væsentlige Henseender korrigeret hans Lære.
Paa flere Punkter (f. Ex. i Spørgsmaalet om den umiddel%
% --- p. 89 ---
\setcounter{footnote}{0}%
\opage{89}bare Opfattelse af Viljen) viser han unægtelig større Klarhed
og Besindighed. Men han er kun i Theorien Pessimist,
eller --- som en af Schopenhauers mere tro Tilhængere,
\emph{Julius Bahnsen}, har sagt --- man véd ikke, om man
skal kalde ham en pessimistisk Optimist eller en optimistisk
Pessimist (Zur Philosophie der Geschichte. Berlin
1872. p. 3). Hartmann anerkjender nemlig den historiske
Udviklings Realitet; den er for ham mere end Isblomster
paa Ruden eller Figurer i Skyerne, for at bruge Schopenhauers
Billeder. Men som Udviklingens Endemaal sætter
han saa Indsigten i Livets Elendighed og en deraf udspringende
Beslutning at standse det. Hvad der hos
Schopenhauer er hver enkelts Opgave, tilfalder hele Slægten
--- og skydes derved ud i en ubestemmelig Fremtid.
I Mellemtiden drive da de optimistiske Illusioner deres
Spil; Enden kommer først, naar de alle have løbet Linen
ud. Hartmann forholder sig til Schopenhauer som de
moderne kristne til den apostoliske Tids troende: disse
ventede Dommedag i en nær Fremtid; for hine ligger den
ude over Horizonten, saa at man foreløbig kan indrette sig
ret hyggelig. Det er derfor karakteristisk, at Hartmann
indbyder sin Læser til at besøge sig for at se, hvor lykkelig
man lever i Pessimistens Hjem.

Hartmann har søgt at oprette, hvad Schopenhauer
forsømte, at gjennemgaa de enkelte Kilder til Nydelse og
undersøge deres Værdi. Men ogsaa han vil i Kraft af sit
nøje Kjendskab til „det ubevidste“ vide bedre Besked om
andres Lykke eller Ulykke end deres egen Følelse. Desuden
er --- som \emph{James Sully} i sin oven anførte Bog har
vist --- hans Undersøgelse højst overfladisk. Især er det
mærkeligt, at han ikke dvæler ved den Lystfølelse, der
udspringer af Virksomhed og Arbejde, men hovedsagelig
holder sig til de passive Nydelser. Tvungne eller af fri
Vilje, i Alvor eller som Leg tilbringe vi jo dog største
% --- p. 90 ---
\setcounter{footnote}{0}%
\opage{90}Delen af vort Liv i Bevægelse og Virksomhed. Dersom
Selvreflexionen gjør ulykkelig, er det jo netop, fordi man
standser paa Vejen, drager Tanken bort fra det Formaal,
den Opgave, hvori den kan finde Hvile. Det sande i
Læren om det ubevidste som Universalmiddel er ikke, at
vi skulle gaa tilbage, men netop at vi skulle gaa frem.
Selvforglemmelsen skal netop tjene til at udvikle det
aandelige Liv videre. Det ubevidste er ikke en Mare,
der tynger os, men en Kilde, hvorfra Livet stedse henter
ny Kraft.

Schopenhauer og Hartmann dele den spekulative
Filosofis Illusion om, at Tanken kan opstille et absolut
Princip. Livet og Erfaringen vise os kun kæmpende
Kræfter, et Spil af Modsætninger, og det er vilkaarligt at
gjøre nogen enkelt af disse til alt. Derimod have disse
Kræfter og Modsætninger en højst forskjellig Værdi for
det menneskelige Liv og dets højere Udvikling. Selv det
mest værdifulde existerer her i Verden kun som en relativ,
kæmpende Magt. Den, der vil vinde det, faar kun
Del i det ved at have Del i dets Kamp for at bestaa og
udvikle sig. Hverken Optimismen eller Pessimismen svarer
til Virkeligheden, naar de skulle gjennemføres konsekvent.
Meliorismen, et af George Elliot (Sully: Pessimism.
p. 399) foreslaaet Udtryk, vil være det rette Ord, naar
man derved forstaar en Tro paa Muligheden af Fremgang
og Udvikling. Og naar vi fastholde den Erfaring, som
Psychologien og Historien ikke mindre end den ydre Natur
indskærper, at Udviklingen gjærne sker ud fra spredte
Punkter, saa fremstiller der sig baade den Opgave og det
Haab, at de i deres Sondring stridende Elementer ved Udviklingens
Lov ville formes og forenes til Helhed og
Harmoni.
\signature{\textbf{\sffamily Harald Høffding.}}

\end{document}
